\chapter{The JML-Inferrer: Heuristic Static Inference}\label{ch:inferrer}

The thesis statement of \cref{ch:introduction} rests on a premise that must be
established before anything else can be built upon it: that specifications can be
recovered from code cheaply enough to be worth having, and that once recovered
they are good enough to be useful. This chapter discharges that premise. It
presents the JML-Inferrer, the single instrument on which the empirical programme
depends, and reports the two results that together answer RQ1: that the inferred
contracts are of demonstrable \emph{fidelity} --- an overwhelming majority are
formally consistent with the code they describe --- and that they deliver
measurable \emph{downstream utility}, improving the unit tests a large language
model writes from them. The conceptual foundations of design by contract, JML, and
extended static checking are laid in \cref{ch:background}; this chapter takes them
as given and concentrates on the instrument and its evaluation.

\section{Architecture}
\label{sec:inf-architecture}

The JML-Inferrer is a static analysis tool, implemented in Java~21 over the
JavaParser abstract syntax tree (AST) library, that reads Java source and emits JML
contracts as annotations inserted in place. The pipeline has three stages. A
parsing stage turns each source file into an AST configured for the Java~21
language level. An inference stage traverses that tree with the visitor pattern,
matching heuristic patterns against syntactic shapes and accumulating candidate
clauses per method. A formatting stage renders the accumulated clauses as JML
comment blocks and writes them back above their methods, preserving the original
source otherwise untouched. Files are processed recursively over a directory tree,
so the tool applies to a codebase exactly as it sits in a repository.

Two design commitments follow the instrument throughout and are worth stating at
the outset because the evaluation depends on them. The first is that inference is
purely static: it requires no test inputs, no instrumentation, and no execution
traces, in deliberate contrast to the dynamic-detector lineage of Daikon, which
generalises likely invariants from observed runs~\citep{ernst2007daikon}. Trading
the dynamic detector's ability to discover numeric relationships from data for the
ability to run on any compilable code without a harness is the exchange that makes
the tool applicable at scale, including to legacy code for which no adequate test
suite exists. The second commitment is determinism: because every heuristic is a
function of the syntax tree alone, the same input always yields the same contract.
This is not merely tidy engineering. It is what allows the downstream study of
\cref{sec:inf-utility} to treat the inferred specification as a controlled variable
rather than a sampled one, and it is the sharpest point of contrast with prompting a
language model to emit JML directly, an approach that returns non-identical
contracts on identical inputs.

\subsection{The contracts the tool infers}
\label{sec:inf-contract-kinds}

The inferrer emits up to five kinds of JML clause, spanning the vocabulary of
design by contract~\citep{meyer1992dbc,leavens2006jml}. \texttt{requires} clauses
state preconditions on the entry state; \texttt{ensures} clauses state
postconditions relating the result and the final state to the entry state;
\texttt{assignable} clauses state the frame, naming the locations a method may
modify; \texttt{loop\_invariant} clauses annotate loops; and class-level invariants
state properties that must hold of an object between method calls. A single method
typically receives a small block combining a guard precondition, one or more
postconditions, and a frame clause; looping methods additionally receive
invariants.

Postconditions carry the subtleties that most exercise the heuristics. A field that
a method modifies must be referred to in its entry-state value under
\lstinline|\old|, so that \lstinline|this.value == \old(this.value) + operand.intValue()|
expresses an increment rather than an impossible fixed point. Branching is
preserved at clause granularity rather than collapsed: a method that returns one
value on one path and another on a second yields two guarded postconditions, each
pairing a return value with the condition that reaches it, of the form
\lstinline|cond ==> \result == e|. This decomposition is not only more faithful; it
is, as \cref{sec:inf-utility} shows, one of the mechanisms by which a downstream
consumer can pick off a contract one clause at a time.

\subsection{Heuristic families organised by WP and SP}
\label{sec:inf-heuristics}

The heuristics are organised by the weakest-precondition and
strongest-postcondition correspondences developed in \cref{ch:background}, which
supply the intuition for what a given syntactic shape licenses. Preconditions are
recovered in the weakest-precondition direction: a guard-and-throw idiom near the
top of a method body --- \lstinline|if (p == null) throw ...|, or a range check that
rejects out-of-bounds indices --- is read as an obligation the caller must have
discharged, and is lifted into a \texttt{requires} clause. Postconditions are
recovered in the strongest-postcondition direction: return statements and field
assignments are propagated forward into \texttt{ensures} clauses, with pre-state
field references wrapped in \lstinline|\old| and path conditions attached where the
value is reached conditionally. A family of loop heuristics --- bound-based
invariants for counter loops, accumulator-pattern detection, monotonicity inference
from update expressions, and bounded symbolic unrolling as a fallback --- supplies
\texttt{loop\_invariant} clauses, emitting a conservative weak invariant only when
all of these fail. Method calls are handled in two passes: the first analyses
methods in dependency order and caches each inferred summary, and the second
propagates those summaries across call sites, resolving calls into the standard
library against a curated specification database and treating calls into
unannotated code as uninterpreted.

Throughout, heuristics that would be unsound under Java's two's-complement integer
semantics are gated by type. A multiplication \lstinline|x * x| or an
\lstinline|Math.abs(x)| does not license \lstinline|\result >= 0| on an
\lstinline|int| or \lstinline|long|, because either can overflow to a negative
value; the non-negativity clause is emitted only for floating-point return types,
where it holds. Deliberately foregoing symbolic execution and full
predicate-transformer computation --- both of which would require an SMT decision
procedure and a complete operational semantics for Java --- is what keeps the
inferrer scalable. Pattern matching trades formal completeness for breadth of
applicability, the established bargain of industrial static analysis, and the
resulting partiality is not hidden: no claim is made that a recovered contract is
complete, only that it is sound, useful, and verifiable.

\section{The verification discipline}
\label{sec:inf-verification}

What keeps heuristic inference honest is that every candidate contract is put to a
verifier. Inferred annotations are checked by OpenJML's Extended Static Checker
(ESC)~\citep{cok2014openjml}, which discharges the resulting proof obligations
through the Z3 SMT solver~\citep{demoura2008z3} without executing the program. A
candidate clause that ESC cannot discharge is treated not as an acceptable
approximation to be shipped but as a defect in the inference to be fixed. This
discipline inverts the usual posture of heuristic tools: rather than emitting
plausible output and leaving validation to the user, the inferrer is developed
against the verifier, and a heuristic earns its place in the repertoire only once
the verifier has discharged the clause it produces on a probe example. A
containerised toolchain makes this reproducible, and an end-to-end regression suite
of several hundred verification tests gates every change to the engine against the
formal oracle.

The thesis is careful never to blur two statuses that this discipline makes
distinct. A \emph{verified} specification is one that ESC has discharged against the
implementation: the checker has established the postcondition, invariant, or frame
clause as a proof obligation from the body under the assumed precondition, or, for a
precondition, has found the clause well-formed and its evaluation defined. A merely
\emph{well-formed} specification is syntactically valid and scope-correct JML that
has not been through, or has not survived, that check. The distinction has teeth:
some inferred clauses involve JML constructs the verifier cannot yet decide, and
these are retained as well-formed oracle information without being counted among the
verified. The remainder of the thesis preserves this boundary wherever the two
might be confused, and \cref{ch:discussion} returns to it as a threat to validity.

A specific obstacle in the verification back-end merited a fix rather than a
work-around. Loop-accumulator postconditions are naturally expressed with the JML
quantifiers \lstinline|\sum|, \lstinline|\product|, and \lstinline|\num_of|, but
the stock OpenJML build reports these as \emph{not yet supported}, which would leave
an entire class of accumulator contracts undischargeable irrespective of their
truth. A forked OpenJML build emits SMT-LIB \lstinline|define-fun-rec| definitions
for these quantifiers, giving the solver a recursive definition to reason with and
unblocking the accumulator-loop arithmetic that would otherwise stall. The fork is
bundled with the toolchain so the validation pipeline remains reproducible. The
strictness configuration is held fixed --- operational integer semantics with
overflow checks in code, mathematical integers in specifications, arithmetic
exceptions treated as failures, and unannotated references treated as nullable ---
so that any change in the discharge outcome is attributable to the inference and not
to a shifting oracle.

\section{Inference fidelity}
\label{sec:inf-fidelity}

The first half of RQ1 asks how faithful the inferred contracts are. Fidelity is
reported on two independent axes measured over a corpus of 312 methods drawn from
eleven classes of Apache Commons Lang, a mature and diverse Java library whose
preponderance of deterministic pure functions suits it to this kind of study. The
first axis is \emph{coverage}: the proportion of methods for which a non-trivial
contract is emitted. The inferrer emits at least one non-trivial precondition or
postcondition for 309 of the 312 methods, a coverage rate of 99.0\%. The three
uncovered methods are overloads that delegate to a standard-library operation for
which no specification-database entry was available, whereupon the engine emits an
empty contract rather than guess --- the conservative limit of interprocedural
summarisation rather than a heuristic failure. For comparison, a documentation-based
baseline that recovers exception preconditions from Javadoc covers under half the
same corpus, and prompting a language model to emit JML covers every method but
with only 72.4\% run-to-run consistency; the static inferrer, deterministic by
construction, is perfectly consistent.

The second axis is \emph{discharge}: the proportion of emitted clauses that ESC can
verify against the implementation. Over all 985 inferred clauses, the checker
discharges 92.9\%, broken down by clause type in \cref{tab:inf-discharge}. Two
readings guard the figure. First, discharge is a soundness measure, not a
completeness one: a vacuous clause would discharge too, so the number should be read
as ``no clause was found inconsistent with the implementation'', which is why the
coverage and structural-strength measures are reported alongside it. Second,
discharge semantics differ by clause type --- for preconditions the clause is
assumed at entry and reported as failing only when ill-formed, so its column is a
well-formedness check rather than a proof of necessity, whereas postconditions,
invariants, and frame clauses are genuine proof obligations. The clauses that do not
discharge separate cleanly into a small \emph{flagged} set, routed to a review queue
and excluded from any downstream use, and an \emph{unsupported} set involving
constructs the solver cannot yet decide. The residual non-discharge concentrates in
a few identifiable patterns --- recursive multiplicative growth, quantified
invariants over mutable arrays, and the way OpenJML translates for-each loops ---
that are properties of the verification technology rather than of the heuristics,
and that delimit what can currently be certified rather than what can be emitted.

\begin{table}[ht]
\centering
\caption{OpenJML ESC discharge over all inferred clauses on the 312-method corpus.
\emph{Discharged} clauses are verified against the implementation; \emph{flagged}
clauses fail discharge and are surfaced for review; \emph{unsupported} clauses
involve constructs the checker cannot yet decide.}
\label{tab:inf-discharge}
\begin{tabular}{lrrrr}
\toprule
\textbf{Clause type} & \textbf{Inferred} & \textbf{Discharged} & \textbf{Flagged} & \textbf{Unsupported} \\
\midrule
Preconditions   & 312 & 296 (94.9\%)  & 11 (3.5\%) &  5 (1.6\%)  \\
Postconditions  & 312 & 268 (85.9\%)  & 22 (7.1\%) & 22 (7.1\%)  \\
Loop invariants &  53 &  43 (81.1\%)  &  4 (7.5\%) &  6 (11.3\%) \\
Assignable      & 308 & 308 (100.0\%) &  0 (0.0\%) &  0 (0.0\%)  \\
\midrule
\textbf{Overall} & \textbf{985} & \textbf{915 (92.9\%)} & \textbf{37 (3.8\%)} & \textbf{33 (3.4\%)} \\
\bottomrule
\end{tabular}
\end{table}

One caveat on fidelity carries through the thesis. The oracle against which clauses
are discharged is the implementation, not the documented intent: ESC confirms that a
clause is consistent with the code it was inferred from, not that it matches the
abstract API contract. For a method whose intended contract is strictly weaker than
its implementation, the inferred clause will be sound against the body yet
over-specified against the interface. This is intrinsic to code-derived inference
and cannot be detected by a discharge oracle that is consistent with the
over-specification by construction; \cref{ch:discussion} records it, and
\cref{ch:compositional} shows how a second analysis pass recovers well-formed
preconditions that single-pass heuristics miss.

\section{Downstream utility: specifications as oracles}
\label{sec:inf-utility}

The second half of RQ1 asks whether these contracts are useful. The concern is not
academic: a specification that is faithful but inert would not vindicate the thesis.
The chosen demonstration is unit-test generation by a large language model, singled
out not because it is the most important use of inferred specifications --- later
chapters pursue compositional verification and agentic integration --- but because
its effect can be measured cleanly. Language models are fluent at the mechanical
part of a test, the prefix that drives a program into a state, but weak at the
oracle, the assertion of what should then be true, because the oracle requires
knowing what the method is \emph{supposed} to do. A signature supplies no such
knowledge and a body supplies too much, tempting the model to encode the
implementation's actual behaviour, bugs included. A specification is precisely the
independent statement of intent that both lack.

The study holds the model fixed and varies only the material supplied alongside an
otherwise identical method signature, over four phases forming an ordered ladder of
information density. P1 supplies the signature alone. P2 adds a fixed,
method-agnostic checklist of edge-case categories --- the control for the rival
hypothesis that any structured guidance, rather than specifications in particular,
raises test quality. P3 adds the inferred JML contract for the method, used without
modification: the treatment condition. P4 substitutes the full source body, an
upper bound available when the model can read the implementation directly. Each
phase was run five times per method to control for sampling non-determinism, and the
generated tests were evaluated by a fully automated pipeline that compiled them, ran
them against the reference implementation, and computed a mutation score under PIT.
Mutation score is the headline metric because it is the available objective proxy
for oracle strength: a suite that kills more mutants is one whose assertions
distinguish more behavioural deviations.

\cref{tab:inf-phases} reports the corpus means. The scores rise monotonically along
the information ladder, from 27.5\% at P1 to 94.8\% at P4. The principal contrast is
P1 against P3: supplying the inferred contract raises the mutation score by 54.3
percentage points, a very large effect. Read against the full-source upper bound,
P3 captures roughly $(81.8 - 27.5)/(94.8 - 27.5) \approx 81\%$ of the achievable
gain from signature-only to full-source generation, and does so within an input an
order of magnitude smaller --- five to twelve lines of annotation against the thirty
to a hundred and fifty lines of body that P4 supplies. The residual gap to P4 is
largely the answer-copying that source code affords and a contract deliberately
withholds: P3 exposes the intended behaviour without the implementation that may
deviate from it.

\begin{table}[ht]
\centering
\caption{Test-generation results by phase, as per-method means across the 312-method
corpus. \emph{Tests} is the average number of generated tests; \emph{Mutation} is
the PIT mutation score. Scores rise monotonically with information density, yet P3
attains a higher mutation score than P2 while generating fewer tests.}
\label{tab:inf-phases}
\begin{tabular}{llrr}
\toprule
\textbf{Phase} & \textbf{Input supplied with signature} & \textbf{Tests} & \textbf{Mutation \%} \\
\midrule
P1 & Signature only                    &  9.0 & 27.5 \\
P2 & Signature + edge-case guidance    & 38.0 & 68.2 \\
P3 & Signature + inferred specification & 33.5 & 81.8 \\
P4 & Full source body                  & 54.7 & 94.8 \\
\bottomrule
\end{tabular}
\end{table}

The most instructive result is the comparison of P3 with the guidance control P2.
The edge-case checklist yields the highest test count of any phase, 38.0 tests per
method, yet its mutation score is 13.6 percentage points \emph{below} P3, which
attains 81.8\% while generating around 12\% fewer tests. The divergence is the
cleanest single piece of evidence for the thesis premise. Three patterns account
for it: the checklist prompts near-duplicate cases that inflate test count without
expanding the set of killed mutants; in the absence of an oracle its tests fall back
on weak assertions such as ``is not null'' or ``does not throw'' that cannot
distinguish a mutant which changes a return value; and P3's tests concentrate on
individual contract clauses, giving a higher mutant-kill density per test. What
distinguishes the treatment from the control is not that the model is told to be
thorough but that it is told \emph{what the method should do}. Test quantity is
therefore a poor proxy for oracle quality, and the benefit being measured is
specifications in particular, not structured prompting in general --- had the study
reported test count alone it would have concluded, wrongly, that generic guidance
was the strongest intervention.

The benefit is uneven across method kinds in a way that confirms the mechanism.
Control-flow methods gain most, because a guarded postcondition supplies an oracle
for each branch that a signature leaves opaque, and utility methods gain substantially
from preconditions that capture their input constraints; accessors gain least, not
because a contract adds nothing but because their signatures already pin down their
behaviour. Within the looping methods, those for which a non-trivial invariant was
inferred outscore those left with a weak invariant, so the downstream benefit tracks
the strength of the inferred specification. The right predictor of gain is the gap a
contract fills over the signature, not its strength in isolation --- and the two ends
of the pipeline are coupled in exactly the direction the thesis predicts: where the
inferrer produces a strong contract the model produces strong tests.

\section{Summary}
\label{sec:inf-summary}

This chapter has presented the JML-Inferrer and answered RQ1 on both axes. On
fidelity, the tool emits a non-trivial contract for 99.0\% of a 312-method corpus,
and OpenJML's Extended Static Checker discharges 92.9\% of the resulting clauses
against the implementation, with the residual failures traced to identifiable limits
of the verification technology rather than to unsound inference. On utility, a
controlled four-phase study shows that supplying the inferred contract to a language
model raises the mutation score of its generated tests by 54.3 percentage points
over the signature-only baseline, capturing roughly four-fifths of the achievable
gain in an order of magnitude smaller input, and --- most tellingly --- outperforming
a generic edge-case checklist that generates more tests but weaker oracles.
Automatically inferred, heuristic, partial specifications are thus both faithful and
useful. That conclusion is the foundation the rest of the thesis builds on:
\cref{ch:compositional} shows how a second, weakest-precondition pass over the call
graph strictly extends this inference and turns callee contracts into a basis for
version-compatibility reasoning, and \cref{ch:integration} places inferred contracts
at the centre of an AI-native development pipeline as the machine-checkable spine its
verification loop requires.
